\documentclass[11pt]{article}
\usepackage[T1]{fontenc}
\usepackage[a4paper,margin=1in]{geometry}
\usepackage{amsmath,amssymb,amsthm}
\usepackage[hidelinks]{hyperref}
\usepackage{microtype}

\newtheorem*{question}{Source Open Problem 7.4}
\newtheorem*{answer}{Separate later answer}

\setlength{\parindent}{0pt}
\setlength{\parskip}{0.55em}

\title{Literature Status: A Weak Tile with\\
Neither a Tiling Complement nor a Spectrum}
\author{Run \texttt{fa\_banach\_001} --- lane 0}
\date{9 August 2026}

\begin{document}
\maketitle

\begin{center}
\fbox{\parbox{0.9\textwidth}{
\textbf{Status: literature already answered.}
Open Problem~7.4 in arXiv:2209.04540 was explicitly answered in the
affirmative by Kiss--Londner--Matolcsi--Somlai, arXiv:2410.04948.
This note records that identification; it claims no new theorem.
}}
\end{center}

\section{Original question}

Kolountzakis, Lev, and Matolcsi's \emph{Spectral sets and weak tiling},
arXiv:2209.04540, asks the following in Section~7.4 on arXiv PDF
page~18~\cite{source}.  Here ``weakly tiles its complement'' means that
there is a positive locally finite measure $\nu$ with
$\mathbf 1_\Omega*\nu=\mathbf 1_{\mathbb R^d\setminus\Omega}$ almost
everywhere.

\begin{question}
Does there exist a bounded measurable set $\Omega\subset\mathbb R^d$
which can weakly tile its complement by translations, but which can neither
tile the space properly by translations nor be a spectral set?
\end{question}

This is the fourth and last open problem in the source paper.

\section{Explicit separate answer}

The abstract and introduction of Kiss, Londner, Matolcsi, and Somlai's
separate paper \emph{A lonely weak tile}, arXiv:2410.04948, explicitly say
that they answer the question raised in the source paper, which they cite as
reference~[19]~\cite{answerpaper}.  On arXiv PDF page~2 they state the target
question and announce an affirmative answer.

\begin{answer}
There is a finite union $C(k)$ of unit cubes in $\mathbb R^5$ which weakly
tiles its complement, but is neither a proper translational tile nor a
spectral set.
\end{answer}

Section~3 first constructs a finite-group set $P_t$ which weakly pd-tiles
the ambient group but is neither spectral nor a tile.  In the final paragraph
on arXiv PDF page~7, the authors transfer the construction to
$C(k)=P(k)+[0,1)^5\subset\mathbb R^5$.  For sufficiently large $k$, standard
lifting results preserve non-tiling and non-spectrality, while the lifted
pd-tiling measure has a unit mass at the origin; removing that mass gives the
required weak tiling of $C(k)^c$.

The provenance is explicit rather than inferred: the later paper names and
cites the source question, and Máté Matolcsi coauthored both papers.

\section{Scope and search status}

This completely answers Open Problem~7.4 in the affirmative.  It does not
answer the other three questions in Section~7: nowhere-dense spectral sets
in dimensions $d\ge2$, spectrality of the second factor in certain product
sets, or whether every extremal weak-tiling measure for a convex body comes
from a proper tiling.

A bounded search through 9 August 2026 checked the run indexes, the exact
wording ``weakly tiles its complement'', and arXiv/literature searches for
weak tiles that are neither spectral nor proper tiles.  The decisive exact
match was arXiv:2410.04948.  No new mathematical result is claimed here.

\begin{thebibliography}{9}
\bibitem{source}
M. N. Kolountzakis, N. Lev, and M. Matolcsi,
\emph{Spectral sets and weak tiling},
Sampl. Theory Signal Process. Data Anal. \textbf{21} (2023), article 31;
arXiv:2209.04540. doi:10.1007/s43670-023-00066-w.

\bibitem{answerpaper}
G. Kiss, I. Londner, M. Matolcsi, and G. Somlai,
\emph{A lonely weak tile},
Expositiones Math. \textbf{44}(1) (2026), article 125636;
arXiv:2410.04948. doi:10.1016/j.exmath.2024.125636.
\end{thebibliography}

\end{document}
